\section{Aplicación móvil}
\label{sec:aplicacion-movil}

\subsection{Arquitectura MVVM}

La aplicación móvil \texttt{apps/mobile} se desarrolló en Kotlin con
Jetpack Compose, siguiendo el patrón arquitectónico Modelo-Vista-VistaModelo
(MVVM) recomendado oficialmente para Android \cite{android2024architecture}.
\texttt{profile}, \texttt{qr}, \texttt{notifications}, \texttt{reservas})
vive en su propio paquete bajo \texttt{features/}, con capas internas
\texttt{presentation} (composables y \texttt{ViewModel}) y \texttt{data}
(repositorios y fuentes remotas/locales), replicando a menor escala la
separación de responsabilidades del backend. La configuración de
\emph{build} fija \texttt{compileSdk}/\texttt{targetSdk} en 34 y
\texttt{minSdk} en 26.

\subsection{Autenticación segura}

El token de sesión se almacena mediante
\texttt{EncryptedAuthStorage.kt}, apoyado en la librería
\texttt{androidx.security:security-crypto}, que cifra el contenido de
\texttt{SharedPreferences} en reposo. El \texttt{AuthViewModel}
gestiona el flujo de inicio de sesión contra el
\texttt{auth-service} y coordina la renovación automática del token
cuando expira, sin requerir intervención del usuario.

\subsection{Módulo de incidentes y notificaciones}

El módulo de incidentes permite registrar y consultar novedades sobre
el estado de los laboratorios. Las notificaciones push se apoyan en
Firebase Cloud Messaging (\texttt{ScliFirebaseMessagingService.kt},
\texttt{NotificationHelper.kt}). Es importante documentar con
honestidad el estado real de esta capacidad: el proyecto todavía no
cuenta con el archivo \texttt{google-services.json} necesario para
activar el complemento de Firebase (el \texttt{build.gradle.kts}
aplica el complemento condicionalmente, solo si el archivo existe), por
lo que el envío de notificaciones push no es funcional en el estado
actual del repositorio, aunque el código de la capacidad está
completo y listo para activarse en cuanto se configure la credencial.

\subsection{Escaneo de código QR}

La segunda capacidad del dispositivo exigida por la guía es el acceso
a la cámara, implementado mediante CameraX
(\texttt{camera-camera2}, \texttt{camera-lifecycle}, \texttt{camera-view})
y ML Kit Barcode Scanning para la detección del código. El flujo
completo, implementado en \texttt{QrScanScreen.kt},
\texttt{QrViewModel.kt} y \texttt{QrRepository.kt}, consiste en
escanear el código QR de un laboratorio, resolver el identificador
UUID capturado, y consultar el endpoint de detalle completo expuesto
por el patrón Facade del backend (\texttt{LaboratorioDetalleFacade}),
el mismo que reutiliza el panel web de monitoreo
(Sección~\ref{sec:aplicacion-web}). La funcionalidad está cubierta por
tres pruebas instrumentadas (\texttt{QrFlowTest.kt}), que verifican el
caso de un QR válido mostrando el detalle del laboratorio, un UUID
inválido mostrando un mensaje de error, y una falla de red mostrando
el mensaje correspondiente del \emph{gateway}.

\begin{figure}[htbp]
\centering
\includegraphics[width=0.7\linewidth]{figuras/qr-scan.png}
\caption{Flujo de escaneo de código QR}
\label{fig:qr-scan}
\end{figure}

\subsection{Módulo de reservas móvil}
% [PENDIENTE -- FREDDY] Describir CalendarioScreen, NuevaReservaScreen,
% ReservasScreen, SolicitudDetalleScreen, con cache local Room para
% modo sin conexion.

\subsection{Capacidades del dispositivo utilizadas}

La Tabla~\ref{tab:capacidades-movil} resume las dos capacidades del
dispositivo exigidas por la guía y su estado real de funcionamiento.

\begin{table}[htbp]
\centering
\caption{Capacidades del dispositivo utilizadas}
\label{tab:capacidades-movil}
\begin{tabular}{lll}
\toprule
Capacidad & Uso & Estado \\
\midrule
Cámara (CameraX + ML Kit) & Escaneo QR de laboratorio & Funcional \\
Notificaciones push (FCM) & Alertas de incidentes/reservas & Pendiente de credencial \\
\bottomrule
\end{tabular}
\end{table}